\documentclass[10pt,letterpaper]{article}
\usepackage[T1]{fontenc}
\usepackage{lmodern,amsmath,amssymb,amsthm,booktabs,array,microtype}
\usepackage[letterpaper,top=0.83in,bottom=0.88in,left=0.94in,right=0.94in,headheight=14pt,headsep=18pt]{geometry}
\usepackage{xcolor,fancyhdr,enumitem,needspace}
\definecolor{ink}{HTML}{16374A}
\definecolor{linkblue}{HTML}{15576F}
\usepackage[colorlinks=true,linkcolor=linkblue,citecolor=linkblue,urlcolor=linkblue,bookmarksnumbered=true]{hyperref}
\hypersetup{pdftitle={Cyclotomic valuations in graph towers and progression-free sets over finite rings},pdfauthor={Research report prepared for Vladimir Reshetnikov}}
\newtheorem{theorem}{Theorem}[section]
\newtheorem{proposition}[theorem]{Proposition}
\newtheorem{lemma}[theorem]{Lemma}
\newtheorem{corollary}[theorem]{Corollary}
\theoremstyle{remark}\newtheorem{remark}[theorem]{Remark}
\setlength{\parindent}{1.15em}
\setlength{\parskip}{1.8pt plus 0.7pt minus 0.3pt}
\setcounter{tocdepth}{1}
\pagestyle{fancy}\fancyhf{}
\lhead{\small\color{ink} Proof methods beyond Dirichlet GRH}
\rhead{\small\color{ink} Cyclotomic and additive methods}
\cfoot{\small\thepage}
\renewcommand{\headrulewidth}{0.25pt}
\setlength{\emergencystretch}{2em}
\raggedbottom
\begin{document}
\thispagestyle{plain}
\noindent{\small\sffamily\bfseries\color{linkblue} RESEARCH REPORT\hfill 25 NOVEMBER 2026}\par
\vspace{.55cm}
\noindent{\LARGE\bfseries\color{ink} Cyclotomic valuations in graph towers}\par
\noindent{\LARGE\bfseries\color{ink} and progression-free sets over finite rings}\par
\vspace{.18cm}
\noindent{\large Two applications of methods in the Lean proof of Dirichlet GRH}\par
\vspace{.30cm}
\noindent\rule{\linewidth}{.5pt}\par
\paragraph{Abstract.}
We determine the full spanning-tree valuation formulas for the five graph
towers in DuBose and Vallières, Section 7, and for all nine rows of
Kataoka's Table 7, including their stated starting levels. A cyclotomic
valuation used internally in the Lean proof reduces the calculation to the
first nonzero term of a Laplacian evaluated at a character. For an abelian
Cayley tower of any rank whose valency is smaller than $p$, a finite sum
of moments over projective space gives $v_p$ of the number of spanning
trees from the first level onward. A separate method in the Lean source,
a Behrend construction on exponents, yields a lower bound for
progression-free subsets of polynomial rings over products of fields.
Combined with a ternary-cube upper bound, it gives a sharp criterion:
positive extremal upper density survives along $\prod_{i\le b}\mathbb F_{q_i}$
as $b\to\infty$ exactly when $\sum_iq_i^{-2}$ is finite.
Each argument is independent of the truth of RH and GRH.
\tableofcontents

\section{What transfers from the Lean source}\label{sec:source}
The readable edition at commit
\href{https://github.com/lsernau/grh-lean/tree/72fead102f663d5c7a407964a13dc0cf9d977207}{\texttt{72fead102f66}}
proves five statements. Its main result gives nonvanishing for Dirichlet
$L(s,\chi)$ with $\Re s>1/2$, taking account of the pole of the principal
character. The critical-strip and primitive-character statements, and
Mathlib's Riemann hypothesis, are explicit consequences in
\texttt{src/Main.lean}. The public theorem concerns Dirichlet characters
of positive integer modulus. The release reports rechecking with Lean
4.35.0, Mathlib revision \texttt{56a00585}, and three independent proof checkers,
with only \texttt{propext}, \texttt{Classical.choice} and \texttt{Quot.sound}.
We inspected the pinned source and checked its public statements and the internal arguments cited below. The formal replays cited here are the
release's own. The proofs of the transfers in this report are mathematical
arguments accompanied by exact computations.

The analytic route compares a short spanning family of integral forms with
a lower bound for the determinant of a chain of theta lifts. The resulting
one-sided inequality controls a weighted sum of prime powers congruent to
$1$ modulo $q$, with a quadratic character separating split and inert
behavior. The transform of the remainder detects orders of zeros. A family
of quadratic twists excludes cancellation at the zero under consideration,
and a one-sided Landau argument completes the reduction. Transferring that
whole step to another $L$-function would require both a suitable
arithmetic bound and a way to control the twisted zeros. The applications
here start at two earlier points in the proof. See the repository guide and
Alpöge's account \cite{LeanRepo,Alpoge}.

\begin{itemize}[nosep,leftmargin=1.5em]
\item In \texttt{Lib/PadicHodge/PeriodRings/AxTheorem.lean}, the source
evaluates the absolute value of $\zeta-1$ for a primitive root of
prime-power order. The proof of \texttt{exists\_norm\_normalizedTrace\_le}
in \texttt{Lib/PadicHodge/TateSen.lean} compares the orders of the terms
of a polynomial in this uniformizer. Their distinct residues modulo the
ramification index prevent cancellation. A related specialization
argument in \texttt{New/Congruences/Selmer/RootsOfUnity.lean} uses the
same radii to choose torsion points for polynomial tests. In a graph
tower we can calculate the order at every character and sum it.
\item In \texttt{New/Reduction/Data/Existence.lean}, the proof of
\texttt{exists\_three\_term\_free} encodes points of a discrete sphere
by digits and rules out weighted three-term congruences between them.
The proof then uses this set to separate certain multiplicative orbits of
places. For a product of polynomial rings, a sum of degrees encodes the
multiplicative progressions. An additional, independent upper bound comes
from density Hales--Jewett \cite{DHJ}.
\end{itemize}

We compared the TeX contents of the 41 earlier ZIPs in
\href{https://github.com/VladimirReshetnikov/ProveIt/tree/5ed40f562710fcc2770bbb463dba528a6727f7ae/docs/incoming}{the pinned ProveIt incoming directory}.
The earlier finite-ring report studied square-zero local algebras. Other
reports developed sharp torsion specialization, including abelian local
coefficient fields. The graph formulas and the reduced-ring density
criterion below did not appear in that collection. We attribute the
numerical graph predictions to their authors and the underlying
cyclotomic and additive tools to their sources. No historical priority
claim is needed for the deductions.

\section{Five numerical patterns in graph towers}\label{sec:graphs}
DuBose and Vallières prove that the valuation of the number of spanning trees
in a $\mathbb Z_p^d$-tower of finite graphs eventually agrees with a polynomial
in $p^n$ and $n$. Section 7 of their paper fits candidate polynomials to finite
data for five two-dimensional towers \cite{DV}. The existence theorem alone
does not specify those polynomials or the first level where they hold.
The following computations determine both for the five examples.

For a list $S$ of integer vectors, use the group $G_n=(\mathbb Z/p^n)^d$ as
vertices. For each $s\in S$ and $v\in G_n$ put an undirected edge
between $v$ and $v+s$. Retain parallel edges, including both edges if a
step has order two. Assume the reductions of $S$ span $\mathbb F_p^d$,
so these graphs are connected. Write $T_n=v_p(\kappa_n)$, with $\kappa_n$
the number of spanning trees and $v_p(p)=1$.

\begin{theorem}\label{thm:five}The five towers of \cite[\S7]{DV}
have the following valuations. Write $e_1=(1,0)$ and $e_2=(0,1)$.
\begin{center}\small
\begin{tabular}{ccll}
\toprule
Example & $p$ & generators $S$ & $T_n$\\
\midrule
1 & 2 & $e_1,e_2$ & $(2n+4)2^n-6n-1$ \\
2 & 2 & $e_1,e_1,e_2,e_2$ & $4^n+(2n+4)2^n-6n-2$ \\
3 & 2 & $(1,5),(0,3),(1,2),(0,1)$ & $(n+33/4)2^n-4n-1$ \\
4 & 3 & $e_1,e_2$ & $4\cdot3^n-2n-4$ \\
5 & 3 & $(1,0),(2,3),(1,1)$ & $20\cdot3^{n-1}-2n-8$ \\
\bottomrule
\end{tabular}
\end{center}
Rows 1, 2, 4 and 5 hold for every $n\ge1$. Row 3 holds for $n\ge4$,
and its first three values are $5,19,65$.
\end{theorem}

\subsection{Reading a character at its conductor}
Fix a primitive $p^h$-th root $\zeta$ in $\overline{\mathbb Q}_p$, set
$e=\varphi(p^h)$ and $t=\zeta-1$. The polynomial $\Phi_{p^h}(1+T)$
is Eisenstein, so $v_p(t)=1/e$. This is the exact radius calculation used
by the repository's torsion specialization argument. The graph application
uses this identity at \emph{every} character of a given conductor.

For $u\in(\mathbb Z/p^h)^d$ primitive (at least one coordinate is a unit),
the Laplacian eigenvalue at that character is
\[
 F_u(\zeta)=\sum_{s\in S}(2-\zeta^{\langle u,s\rangle}
                                  -\zeta^{-\langle u,s\rangle}).
\]
Fourier characters diagonalize translations. Over the complex numbers all
these eigenvalues are positive for nontrivial characters by connectedness.
The matrix-tree theorem gives the product of nonzero eigenvalues divided by
$p^{dn}$. Embedding the algebraic identity into $\overline{\mathbb Q}_p$
and grouping by conductor yields
\begin{equation}\label{eq:shells}
 T_n=-dn+\sum_{h=1}^n W_h,\qquad
 W_h=\sum_{[u]} e\,v_p(F_u(\zeta)).
\end{equation}
The brackets denote orbits of primitive vectors under multiplication by units
modulo $p^h$. The action is free and its members are cyclotomic conjugates,
all of the same valuation. This orbit decomposition is already used by
DuBose and Vallières. Our evaluations of the individual orders give the
all-level formulas.

Two elementary rules will be useful. If an integral power series in $T$
has first unit coefficient at degree $k$, and $k<e$, then its value at
$t$ has valuation $k/e$. The same conclusion holds for $k<e+2$ if the
constant and linear coefficients vanish identically. Both assertions follow
by comparing the valuations term by term: lower coefficients divisible by
$p$ contribute above the candidate, and higher powers do also. For arbitrary
$k$ one can reduce the polynomial modulo $\Phi_{p^h}(1+T)$, obtaining
$\sum_{j<e}c_jT^j$; the normalized order is exactly
$\min_{c_j\ne0}\{e\,v_p(c_j)+j\}$. These integers lie in distinct residue
classes modulo $e$. This gives the exact small-conductor check used below.

\subsection{The dyadic square torus and its doubled edges}
Write $X=\zeta^a$, $Y=\zeta^b$ and $F=4-X-X^{-1}-Y-Y^{-1}$.
The identity
\[
 F=4-(X+Y)(1+(XY)^{-1})
\]
makes the exceptional characters explicit. For $h\ge3$, put $e=2^{h-1}$.
For a mixed-parity primitive pair, both $a-b$ and $a+b$ are odd.
Each factor in the product therefore has valuation $1/e$, giving normalized
order $2$. There are $2e$ mixed-parity orbits, contributing $4e$.

For the remaining orbits normalize $a=1$, so $b$ runs through the $e$ odd
residues modulo $2^h$. The normalized orders of the factors are those of
$1+\zeta^{1-b}$ and $1+\zeta^{1+b}$. For an exponent $k$ put $r=k+e$
modulo $2^h$. If $r\ne0$, this order is $2^{v_2(r)}$. If $r=0$ the factor
vanishes. For either sign, $1\pm b+e$ runs exactly once through the even
residues. Of a given pair of exponents, one is $2$ modulo $4$, so its order
is $2$. The other, when finite, has order at most $e$.
Thus a finite product has normalized order at most $e+2<2e$, and subtracting
from $4$ preserves that order. At a vanishing factor the order of $F$ is
$2e$ instead.

For either sign, the sum over the nonzero even residues is
$\sum_{i=1}^{h-1}(e/2^i)2^i=(h-1)e$. At each of the two exceptional orbits
the other factor has already contributed $2$ to this sum, so each requires a
correction of $2e-2$. The odd--odd contribution is
$2(h-1)e+2(2e-2)$. Consequently
\[
 W_h=(h+3)2^h-4\quad(h\ge3).
\]
At $h=1$ the eigenvalues are $4,4,8$, giving $W_1=7$.
At $h=2$ the four mixed orbits each contribute $2$, and the two odd--odd
orbits each contribute $4$, giving $W_2=16$. Equation~\eqref{eq:shells}
gives row 1 by summation. Doubling each edge multiplies the number of
spanning trees by $2^{4^n-1}$, since a tree uses $4^n-1$ edges with two
independent parallel choices each. This proves row 2.

\subsection{The four-loop dyadic example with unequal slopes}
In Example~(3) of Section~7 of DuBose and Valli\`eres, the four directions are $(1,5),(0,3),(1,2),(0,1)$. Write $T_n=v_2(\kappa_n)$ for its spanning-tree valuation. Then
\[
 T_n=n2^n+\tfrac{33}{4}2^n-4n-1\qquad(n\ge4),
 \qquad (T_1,T_2,T_3)=(5,19,65).
\]
To see this, write $F=8-\sum_{(u,v)}(X^uY^v+X^{-u}Y^{-v})$.
At conductor $2^s$, take $\zeta$ primitive, $e=2^{s-1}$, and let $W_s$ be the sum of $v_2(F(\zeta^a,\zeta^b))$ over pairs modulo $2^s$ with at least one odd entry. The two vectors $(1,2),(0,1)$ generate the group, so the only zero of $F$ among these characters is the excluded trivial character. The matrix-tree theorem and character diagonalization give $T_s-T_{s-1}=W_s-2$.

The decisive local computation is the following identity over $\mathbb F_2$:
\[
 XY^5\overline F=(Y+1)\big[(Y^2+Y+1)(X^2Y^7+1)+XY^2(Y+1)^5\big].
\]
Substitute $X=(1+z)^a$, $Y=(1+z)^b$, using nonnegative representatives of $a,b$. For $b\ne0$, put $B=2^{v_2(b)}$, $C=2^{v_2(2a+7b)}$. In characteristic two, $(1+z)^j-1$ has order $2^{v_2(j)}$. Since $C$ is a power of two, $C\ne5B$. The resulting order of $\overline F$ is exactly $D=B+\min(C,5B)$. Whenever $D<e$ this gives $ev_2(F(\zeta^a,\zeta^b))=D$: in the integral expansion at $z=\zeta-1$ the lower powers have even coefficients, the coefficient of $z^D$ is odd, and $v_2(\zeta-1)=1/e$. One can use the polynomial $XY^5 F$ to avoid negative exponents throughout.

For $s\ge5$, the sums over the following disjoint classes exhaust the conductor shell.
\begin{center}
\begin{tabular}{lll}
\hline
Class & Normalized order $ev_2(F)$ & Contribution to $W_s$\\
\hline
$b$ odd & $2$ & $4e$\\
$b\equiv2\pmod4$ & $6,10,12$ & $17e/4$\\
$v_2(b)=k$, $2\le k\le s-2$ & $2^k+2$ & $(s-3)e+e-4$ (all $k$)\\
$b=0$ & $e+2$ & $e+2$\\
$b=e$ & $2e$ & $2e$\\
\hline
\end{tabular}
\end{center}
Indeed, for fixed $b\equiv2\pmod4$, as $a$ runs through odd residues, $C$ is $4,8,$ or at least $16$ in proportions $1/2,1/4,1/4$. For $k\ge2$ we have $C=2$ and $e/2^k$ eligible values of $b$ at depth $k$. All three generic rows have $D<e$. For the two exceptional rows, the integer identities $F(X,1)=2(2-X-X^{-1})$ and $F(X,-1)=12$ give the listed orders directly. Summation yields
\[
 W_s=(s+37/4)2^{s-1}-2\quad(s\ge5).
\]
The small-shell totals $W_1,W_2,W_3,W_4$ are $7,16,48,116$. For an exact short verification, use one representative per Galois orbit, $(a,1)$ with arbitrary $a$ and $(1,b)$ with $b$ even. Reduce $F(Z^a,Z^b)$ modulo $Z^e+1$ and write the result as $\sum_{j<e}c_j(Z-1)^j$. The orbit contributes $\min_{c_j\ne0}(ev_2(c_j)+j)$, since its size is $e$ and these orders are distinct modulo $e$. At $s=1,2,3,4$, the multisets of orbit contributions, using superscripts for multiplicities, are respectively
\[
 \{2^{[2]},3\},\quad \{2^{[4]},4^{[2]}\},\quad
 \{2^{[8]},6^{[2]},8,12\},\quad
 \{2^{[16]},6^{[4]},10^{[2]},16,24\}.
\]
Thus $T_4=179$. Summing the shell recurrence proves the formula for all subsequent levels.


\subsection{The two ternary towers}
For the square torus $p=3$, the coefficient of $T^2$ in
$F_u(1+T)$ is $-(a^2+b^2)$ and the lower two coefficients are exactly zero.
Every primitive pair has $a^2+b^2\ne0$ modulo three. Thus each orbit has
order $2$, even at $h=1$ by the strengthened comparison rule. The number
of orbits is $4\cdot3^{h-1}$, giving row 4 of Theorem~\ref{thm:five}.

For row 5 the projected steps are $a,,2a+3b,,a+b$. At $h=1$ the
eigenvalue is three times the number of nonzero projected steps modulo
three. The four orbits have normalized orders $4,4,2,2$, so $W_1=12$.
For $h\ge2$, classify the primitive direction modulo nine. The following
table gives the first nonzero degree $k$ of $F_u(1+T)$ modulo $(3,T^9)$.
\begin{center}\small
\begin{tabular}{c|rrrrrrrrr|ccc}
$u$ & $(1,0)$ & $(1,1)$ & $(1,2)$ & $(1,3)$ & $(1,4)$ & $(1,5)$ &
$(1,6)$ & $(1,7)$ & $(1,8)$ & $(0,1)$ & $(3,1)$ & $(6,1)$\\
\hline
$k$&4&6&2&6&4&2&4&4&2&2&2&2
\end{tabular}
\end{center}
For an explicit verification put $J=\{a,2a+3b,a+b\}$ as a multiset.
For $1\le j\le6$ the coefficient is the reduction modulo three of
$-\sum_{k'\in J}(\binom{[k']_9}{j}+\binom{[-k']_9}{j})$,
where the representatives lie in $[0,8]$. The constant is zero before reduction.
The table persists at every higher lift since
$(1+T)^9\equiv1\pmod{3,T^9}$. Normalizing a unit multiple amounts to a
change of variable with invertible linear term. In the original Laurent
expansion both the constant and the linear terms vanish exactly; thus
$k\le6<e+2$ makes each displayed $k$ the normalized valuation at every
$h\ge2$. Every displayed orbit has $3^{h-2}$ orbit lifts. The sum of the
table is $40$, so $W_h=40\cdot3^{h-2}$. Summing from $W_1=12$ proves row 5.


\section{Nine suggested formulas of Kataoka}\label{sec:kataoka}
Kataoka's 2026 preprint \cite[\S6.3, Tables 6--7]{Kat} treats dyadic
towers with prescribed leading coefficient $\lambda$. Table 7 lists data
through level six and suggests a full formula for each $0\le\lambda\le8$.
The author verifies the coefficient of $n2^n$ there and leaves the
remaining coefficients unconfirmed. The following result establishes those
coefficients and the indicated ranges. Here $X,Y$ denote voltage
monomials for the two basis vectors, so $X^aY^b$ corresponds to the step
$(a,b)$. A superscript on $X$ is an exponent; a repeated entry means
parallel generators.

\Needspace{15\baselineskip}
\begin{theorem}\label{thm:kataoka}
Put $N=2^n$. For the nine step lists below, the valuation $T_n$ equals the
given expression for every $n$ in the displayed range.
\begin{center}\small
\begin{tabular}{cllc}
\toprule
$\lambda$ & Generators & $T_n$ & First $n$\\
\midrule
0 & $X^3Y^2,X,Y$ & $9N-2n-3$ & 3\\
1 & $X^2Y^2,X,Y$ & $nN+(15/2)N-4n-7$ & 3\\
2 & $X,Y,Y$ & $2nN+2n+1$ & 2\\
3 & $XY,X,Y$ & $3nN+9N-8n-9$ & 1\\
4 & $X^2,X,X,Y,Y$ & $4nN-4N+6n+8$ & 2\\
5 & $X^2Y^2,X^2Y,XY^2,X,Y$ & $5nN+5N-16n+22$ & 3\\
6 & $X^3,X,Y,Y$ & $6nN-10N+14n+14$ & 3\\
7 & $X^3Y^3,X^3Y^2,X^2Y^3,X^2Y,XY^2,XY$ & $7nN+4N-20n+14$ & 3\\
8 & $X^4,X,X,Y,Y$ & $8nN-16N+14n+6$ & 2\\
\bottomrule
\end{tabular}
\end{center}
\end{theorem}

\subsection{Bounded orders in the row $\lambda=0$}
For $u=(a,b)$, the projected steps are $3a+2b,a,b$. Reduce the expansion
of $F_u(1+t)$ modulo $(2,t^8)$, exactly as for the ternary table above.
For $(1,b)$, $0\le b\le7$, the orders are
$6,2,4,2,6,2,4,2$ in sequence. For $(a,1)$, $a=0,2,4,6$, they are all $2$.
These orders are stable on lifting each class modulo eight, and they are
strictly smaller than $e=2^{h-1}$ for $h\ge4$.
The sum of the twelve entries is $36$, so $W_h=36\cdot2^{h-3}$.
The first three exact shell sums $7,18,44$ give $T_3=63$. The recurrence
\eqref{eq:shells} yields $T_n=9\cdot2^n-2n-3$ for $n\ge3$.

\paragraph{Kataoka's row $\lambda=1$.}
Here the three steps are $(2,2),(1,0),(0,1)$. Put
$X=\zeta^a$, $Y=\zeta^b$, $Z=XY$, and write $\nu=e\,v_2$ at conductor
$2^h$, where $e=2^{h-1}$. The Laplacian eigenvalue has the exact factorization
\[
 F=4-(Z+1)H,\qquad
 H=(X+Y)/Z + (Z-1)^2(Z+1)/Z^2 .
\]
Indeed the third pair of translated terms adds $-(Z-Z^{-1})^2$ to the
square-torus eigenvalue. Work first with $h\ge4$.
For mixed parities the square-torus eigenvalue has order $2$ and
$\nu((Z-Z^{-1})^2)=4$. The $2e$ mixed orbits therefore contribute $4e$.

For odd $a,b$ take $a=1$. If $b\equiv-1\pmod4$, the first term of
$H$ has order $2$, since $\nu(X+Y)=\nu(1-\zeta^{b-1+e})=2$.
The second has order at least $12$ or vanishes. Thus $\nu(H)=2$.
Over these $e/2$ orbits, the exponents $1+b+e$ range over all multiples
of four modulo $2^h$. For a nonzero such exponent, say $k$, the factor
$Z+1$ has order $2^{v_2(k)}\le e$. The finite orders sum to
$\sum_{j=2}^{h-1}(e/2^j)2^j=(h-2)e$. One orbit has $Z=-1$, giving
$\nu(F)=2e$. On the others $\nu(F)=\nu(Z+1)+2<2e$. These orbits
together contribute $(h+1)e-2$.

For $b\equiv1\pmod4$ both $Z-1$ and $Z+1$ have order $2$.
Consequently the second term of $H$ has order $6$.
The first has order $4$ for $b\equiv5\pmod8$ and at least $8$
(including infinity) for $b\equiv1\pmod8$, since $8\mid e$.
The two contributions to $H$ have distinct orders in both cases.
The resulting orders of $F$ are $6$ and $8$, each occurring $e/4$
times. Their sum is $7e/2$. We have proved
\[
          W_h=(h+17/2)2^{h-1}-2\qquad(h\ge4).
\]
Direct cyclotomic reduction at $h=1,2,3$ gives the respective shell sums
$7,16,48$. Thus $T_1=5$, $T_2=19$, $T_3=65$. The recurrence
$T_h-T_{h-1}=W_h-2$ now gives
\[
             T_n=(n+15/2)2^n-4n-7\qquad(n\ge3).
\]

\subsection{The four even positive indices and the triangular lattice}
Let $A(U)=2-U-U^{-1}$. For $s\ge1$, put $e=2^{s-1}$ and take $\zeta$
of exact order $2^s$. The conductor shell consists of the pairs
$(X,Y)=(\zeta^a,\zeta^b)$ with at least one exponent odd modulo $2^s$.
If $W_s$ sums the valuations of the Laplacian eigenvalue $F$ over that
shell, the matrix-tree formula gives $T_s-T_{s-1}=W_s-2$.

For $\lambda=2m$, $1\le m\le4$, the eigenvalue is
$F=P_m(A(X))+2A(Y)$, where the integral polynomials satisfy
$P_m(0)=0$ and $P_m(z)\equiv z^m\pmod2$. Explicitly, the four
polynomials are $z$, $z(6-z)$, $z((3-z)^2+1)$, and
$z(4-z)[4-z(4-z)]+2z$. These identities follow from
$A(X^2)=A(X)(4-A(X))$ and $A(X^3)=A(X)(3-A(X))^2$.

Put $L=2^j$, the smallest power of two at least $m$.
Take $s\ge\max(3,j+2)$. If $a\ne0$ then
$t=v_2(A(X))=2\cdot 2^{v_2(a)}/e$. When $t\le1/L$, the odd
coefficient of $z^m$ gives $v_2(P_m(A(X)))=mt\le1$, since lower
coefficients are even and the constant term vanishes.
For odd $a$ this applies regardless of $b$; the $2e^2$ choices
contribute $4me$ to $W_s$. With $a$ even and $b$ odd, each stratum
$v_2(a)=k$ for $1\le k\le s-2-j$ contributes $2me$. There are
$e/2^k$ choices of $a$ and $e$ choices of $b$ in that stratum.
In the remaining $2L$ choices of $a$, either $a=0$, or $t\ge2/L$.
For the latter, the leading monomial has valuation
$mt\ge 2m/L>1+2/e$, and the terms of smaller positive degree
have valuation at least $1+t>1+2/e$. Higher terms also exceed this
bound. The strict bound for $2m/L$ follows at once for
$(m,L)=(1,1),(2,2),(3,4),(4,4)$ and the stated range of $s$.
The valuation of $F$ for all these $a$ is thus $1+2/e$.
In total,
\[
 W_s=\big[2m(s-j)+2L\big]e+4L.
\]

For $\lambda=3$ use the exact identity
$XYF=8XY-(1+X)(1+Y)(1+XY)$. For $s\ge2$, exactly one of $a,b,a+b$
is even. Label this residue $t$. The two odd arguments each give
$v_2(1+\zeta^{\mathrm{odd}})=1/e$. For each choice of the even position
and its value there are $e$ pairs. When $t=e$ the product is zero,
hence $v_2(F)=3$. For all the other even $t$ we have
$v_2(1+\zeta^t)\le1$, hence $v_2(F)=2/e+v_2(1+\zeta^t)<3$.
The sum of $v_2(1+\zeta^t)$ over those remaining even $t$ is $s-1$,
since $-\zeta^t$ takes all nonidentity values in $\mu_e$ once and
$\prod_{\xi\in\mu_e,\,\xi\ne1}(1-\xi)=e$. Thus
$W_s=3(s+4)e-6$.

For clarity, in the order $\lambda=2,4,6,8,3$, the shell totals at
$s=1,2,3$ are $(7,18,36)$, $(10,30,56)$, $(10,28,88)$,
$(10,28,80)$ and $(9,30,78)$, respectively. They can be obtained by
reducing each orbit's eigenvalue modulo $Z^e+1$ and taking the smallest
valuation in the basis $1,(Z-1),\ldots,(Z-1)^{e-1}$. The accompanying
integer checker performs these small computations independently.
Substituting in $T_0=0$, $T_s=T_{s-1}+W_s-2$, then using the shell
formulas, gives the corresponding five rows of Theorem~\ref{thm:kataoka}.

\subsection{Kataoka's two odd rows with a repeated diagonal factor}

The first list contains a basis, and the second contains $(2,1),(1,1)$, another basis. Thus all layers are connected.

Put $d=\lambda-2$, $m=2^s$, $e=2^{s-1}$ and let $F$ be the Laplacian symbol. Direct multiplication gives
\[
 \overline F\doteq (X+1)(Y+1)(XY+1)^d
 \quad\text{in }\mathbb F_2[X^{\pm1},Y^{\pm1}],
\]
where $\doteq$ means equality up to a Laurent monomial. For each primitive pair $(a,b)$ modulo $m$, exactly one of $a,b,a+b$ is even. Write $c$ for this even member and $W=\zeta^c$, $\zeta$ primitive of order $m$. Let $2^t$ be the order of $W$, including $t=0$ for $W=1$. For $t\ge1$ the order of $(1+z)^c-1$ in characteristic two is $B=e/2^{t-1}$, taking $c$ to be its residue in $[0,m)$. Replacing an exponent by a congruent one does not change the formal order below $m$. At the axis $a$ even or $b$ even the contact order is $B+d+1$. At the diagonal $a+b$ even it is $2+dB$. If the contact order is strictly below $e$, it equals the normalized valuation $e v_2(F(\zeta^a,\zeta^b))$ by the uniformizer test. For $s\ge5$ this settles both axes when $t\ge2$, and the diagonal when $t\ge L+1$, where $L=2$ for $d=3$ and $L=3$ for $d=5$.

For the remaining classes let $h=(\lambda-1)/2$ and let $Z$ denote a primitive root of order $m$. In an axis class expand $W^hZ^hF(W,Z)$ as $\sum_j c_j(Z-1)^j$; symmetry treats the second axis. In a diagonal class use $W^h ZF(Z,W/Z)$ instead. The following table gives the smallest pair $(v_2(c_j),j)$ in lexicographic order. A dash means that the generic calculation applies.
\[
\begin{array}{c|cc|cc}
&\multicolumn{2}{c|}{\text{axis}}&\multicolumn{2}{c}{\text{diagonal}}\\
 t&\lambda=5&\lambda=7&\lambda=5&\lambda=7\\\hline
 0&(1,2)&(1,4)&(2,2)&(2,2)\\
 1&(1,2)&(1,4)&(3,0)&(4,0)\\
 2&-&-&(2,0)&(3/2,2)\\
 3&-&-&-&(2,0)
\end{array}
\]
To verify the entries, at $W=1,-1$ respectively the axis coefficient lists, in increasing $j$, are
\[
\begin{array}{c|ll}
5&(0,0,-10,-10,-2)&(8,16,6,-2,0)\\
7&(0,0,-28,-56,-42,-14,-2)&(-16,-48,-56,-32,-10,-2,0).
\end{array}
\]
For the diagonal the lists in increasing $t$ are
\[
\begin{array}{c|l}
5&(0,0,-4),\quad(8,8,0),\quad(-12,-12,0)\\
7&(0,0,-4),\quad(-16,-16,0),\quad
 (-16W,-16W-4,-2W-2),\quad(12W^3,12W^3,0).
\end{array}
\]
These expansions use only the identity $\Phi_{2^t}(W)=0$ and $\Phi_1(W)=W-1$ when $t=0$. For $W$ of order four we have $v_2(1+W)=1/2$, explaining the fractional entry. If the table entry is $(H,J)$, the normalized valuation is $eH+J$: the coefficient valuations lie in $\varphi(2^t)^{-1}\mathbb Z$. For $s\ge5$, the polynomial degree (at most six on the axes and two on the diagonal) is smaller than $e/\varphi(2^t)$ in every tabulated case. Thus the minimum term is unique.

To sum each class over its primitive pairs, for each fixed $W$ there are $e$ choices for the primitive coordinate $Z$. A normalized order $D$ therefore contributes $D$ for each such $W$, and there are $\varphi(2^t)$ choices of exact order $2^t$. In a single axis class the generic sum is $(s-2)e+(d+1)(e-2)$, while its two exceptional classes contribute $2(e+d-1)$. Each axis total is consequently $(s+d+1)e-4$. The generic diagonal sum is
\[
 d(s-L-1)e+2(e-2^L).
\]
Its exceptional classes contribute $9e+2$ for $d=3$ and $17e+6$ for $d=5$. Adding all three classes yields
\[
 W_s=(5s+10)e-14\quad (\lambda=5),\qquad
 W_s=(7s+11)e-18\quad (\lambda=7),\qquad s\ge5.
\]
The four initial valuations of the trees are $(8,34,134,358)$ and $(10,50,154,446)$ respectively. They can be verified by the finite cyclotomic basis calculation given above; the accompanying script recomputes all twelve initial layers. Since $T_s-T_{s-1}=W_s-2$, summation proves the asserted formulas.


\section{A finite-field formula for abelian Cayley towers}\label{sec:moments}
The preceding argument has a uniform form when the prime exceeds the
valency. Fix $m$ steps $s_i$ spanning $\mathbb F_p^d$ modulo $p$, with
$p>2m$. For a projective covector $L\in\mathbb P^{d-1}(\mathbb F_p)$
represented by $u$, define
\[
 r(L)=\min\{r\ge1:\textstyle\sum_{i=1}^m\langle u,s_i\rangle^{2r}
                       \ne0\ \text{in }\mathbb F_p\},
 \qquad A=2\sum_L r(L).
\]
Scaling $u$ preserves the first nonvanishing degree.

\begin{theorem}\label{thm:moments}
For every $n\ge1$,
\[
 \boxed{\displaystyle T_n=A\sum_{h=0}^{n-1}p^{(d-1)h}-dn.}
\]
In particular, the reductions of the steps modulo $p$ determine the
valuation at \emph{every} level. Each $r(L)$ is at most $m$.
\end{theorem}

\begin{proof}
Put $x_i=\langle u,s_i\rangle^2$ in $\mathbb F_p$. If their first $m$
power sums vanished, Newton's identities (with $1,\ldots,m$ invertible)
would give $\prod_i(Z-x_i)=Z^m$. This forces all projections zero and
contradicts spanning. Next take the truncated exponential coordinate
$1+T=\exp z$ modulo $z^{2m+1}$ over $\mathbb F_p$. All denominators used
are units, and the linear term is invertible. The eigenvalue becomes
\[
 -2\sum_{j=1}^m\frac{z^{2j}}{(2j)!}
                 \sum_i\langle u,s_i\rangle^{2j}
                 \pmod {z^{2m+1}} .
\]
Its first nonzero degree is $2r(L)$. The integral series has zero constant
and linear coefficients. Since $2r(L)\le2m\le p-1\le e$, the second
comparison rule proves that the orbit contributes exactly $2r(L)$.
There are $p^{(d-1)(h-1)}$ lifts of each projective residue direction to
conductor $p^h$: normalize the first unit coordinate to one and lift the
other $d-1$ coordinates freely. Formula~\eqref{eq:shells} now gives the claim.
\end{proof}

For example, for the square torus and $p\ge5$, the first moment is
$a^2+b^2$. For $p\equiv3\pmod4$ it never vanishes on a projective point.
For $p\equiv1\pmod4$ its two isotropic directions have nonzero fourth
moment $2a^4$. Consequently $A=2(p+1)$ and $A=2(p+3)$, respectively.
At $p=5$ this gives $T_n=4(5^n-1)-2n$. In rank one, Dion, Lei, Ray and
Vallières had already used the even moments and analyzed the distribution
of the one-variable Iwasawa invariant, including the two-loop case
\cite[Theorem 4.13 and Proposition 4.19]{DLRV}. The formula above counts
all residual projective directions and their lifts in arbitrary rank.

\paragraph{Sharpness in rank two.}
Suppose $m\ge2$ and $2m\mid p-1$. Take $g$ of order $2m$ modulo $p$
and lift the $m$ steps $(g^j,g^{-j})$, $0\le j<m$, to integers.
The first two span. In the direction $(a,b)$ the quadratic moment is
$2mab$, since the sums of $g^{2j}$ and $g^{-2j}$ vanish. Each coordinate
axis has all even moments through order $2m-2$ zero, and the $2m$-th moment is
nonzero. Thus $A=2(p-1)+4m$. This is the largest possible $A$ for
$m$ steps in rank two when their quadratic moment is not the zero polynomial:
a nonzero binary quadratic has at most two projective roots, each of
which can contribute at most $2m$. This proves sharpness for the stated
congruence class of primes. For example $p=13,m=6$ gives $A=48$.

\section{Geometric progressions over products of finite fields}
\label{sec:products}
Asada, Fourakis, Manski, McNew, Miller and Moreland ask in Question~6.5
of their \href{https://doi.org/10.1016/j.ffa.2016.10.002}{2017 paper}
how their density problem changes over other finite coefficient rings.
Their Remark~1.1 specifies polynomial ratios that are nonunits.
A finite product of fields gives a new feature: multiplication by $x$ can be
switched on independently in any subset of the factors. This yields a sharp
criterion for whether a positive upper density can survive as factors are added.

For $R_b=\prod_{i=1}^b\mathbb F_{q_i}$ take
$B_n(R_b)=\{f\in R_b[x]:\text{all coefficients above degree }n\text{ vanish}\}$.
It is the product of the coordinate balls and has size $|R_b|^{n+1}$.
Write $\Delta(R_b)$ for the supremum of
$\limsup_{n\to\infty}|E\cap B_n(R_b)|/|B_n(R_b)|$ over sets $E$
avoiding triples $f,rf,r^2f$ with a nonunit $r\in R_b[x]$.
All constructions below have every component nonzero. The upper bound
uses only progressions with three distinct terms. Thus possible conventions
about repeated terms or zero leave the results unchanged: the proportion of
tuples having a zero component is at most $\sum_i q_i^{-(n+1)}$ at fixed $b$.

\begin{theorem}[Criterion for growing products]\label{thm:products}
For an arbitrary sequence of prime powers $q_1,q_2,\ldots$, the limit
$\lim_{b\to\infty}\Delta(R_b)$ exists. It is positive precisely when
$\sum_i q_i^{-2}$ converges. For every finite $b$,
\[
 \prod_{i=1}^b(1-q_i^{-2})\ \le\ \Delta(R_b)
 \ \le\ \mathbb E\, a_K,\qquad
 K=\sum_{i=1}^b X_i,\qquad
 \Pr(X_i=1)=\frac{3}{q_i^2+q_i+1},
\]
where $X_i$ are independent Bernoulli variables and $a_k$ is the largest
proportion of the ternary word cube $\{0,1,2\}^k$ containing no
combinatorial line, with $a_0=1$. In particular $\Delta(\mathbb F_q^b)\to0$
for fixed $q$. There is an absolute $C>0$ such that for every $b\ge1$
and all $q_i$,
\[
 \Delta(R_b)\ge \exp\!\big(-C\sqrt{\log(2+b)}\big).
\]
\end{theorem}

\subsection{Packing the polynomial ball with word cubes}
A combinatorial line has three words, with a nonempty set of variable
positions reading respectively $0,1,2$ in unison and all other positions
fixed. Fix a field of order $q$. Each nonzero polynomial in its ball
has a unique expression $x^eg$ where $g(0)\ne0$. For each fixed $g$,
partition the exponent interval $0\le e\le n-\deg g$ into disjoint
successive triples starting at $0,3,6,\ldots$, leaving up to two singleton
cells. Put zero in a singleton. For $3j+2\le n$, the number of triples
starting at exponent $3j$ equals $(q-1)q^{n-2-3j}$, the number of choices
of $g$ having nonzero constant term and degree at most $n-2-3j$.
Consequently the fraction of the ball lying in triple cells is
\[
 \theta_{q,n}=3(q-1)\sum_{0\le j\le(n-2)/3}q^{-3j-3}
 \ \longrightarrow\ \frac{3}{q^2+q+1}.
\]

Choose a point uniformly in the product ball, then condition on the cell
of each coordinate. The coordinates in triple cells are independent
uniform letters. Their count $K_n$ has the law of a sum of independent
Bernoulli variables with parameters $\theta_{q_i,n}$. Within a cell
with $k$ active coordinates, each combinatorial line is
$f,rf,r^2f$: take $r_i=x$ on its variable positions and $r_i=1$
elsewhere. The three tuples are distinct, since the variable positions
belong to nonzero chains. Conditional occupancy by any admissible set
is at most $a_k$. Averaging gives an upper bound $\mathbb E a_{K_n}$
even for the finite-ball maximum. At fixed $b$, passage to the limit in
this finite distribution gives the asserted upper bound.

The \href{https://annals.math.princeton.edu/2012/175-3/p06}{density
Hales--Jewett theorem} gives $a_k\to0$. Write
$\mu_b=\sum_{i\le b}3/(q_i^2+q_i+1)$. For $q_i\ge2$ the summands
are between $(12/7)q_i^{-2}$ and $3q_i^{-2}$. Thus $\mu_b\to\infty$
when $\sum_iq_i^{-2}$ diverges. Since $\mathrm{Var}(K)\le\mu_b$,
the probability $\Pr(K<\mu_b/2)$ is at most $4/\mu_b$. The contribution
of larger $K$ is bounded by $\sup_{k\ge\mu_b/2}a_k\to0$. This proves the
vanishing assertion, using only the $2^b-1$ possible nonconstant ratios
whose coordinates are in $\{1,x\}$.

\subsection{Passing from a degree cutoff to a single infinite set}
Let $M_b(n)$ be the largest admissible fraction of $B_n(R_b)$.
Then
\begin{equation}\label{eq:glueproduct}
 \Delta(R_b)=\limsup_{n\to\infty} M_b(n).
\end{equation}
For the nontrivial direction choose $n_j$ so that $M_b(n_j)$ approaches
its upper limit, $n_j-j>2n_{j-1}$, and $n_j-2j>n_{j-1}$.
In a maximizing set at cutoff $n_j$ keep only tuples with all degrees
between $n_j-j$ and $n_j$, taking $n_1>1$. This deletes a proportion
at most $\sum_i q_i^{-j-1}$. Form the union of the retained boxes.
In any progression in this union all $f_i,r_i$ are nonzero, and
the three degrees in each coordinate are nondecreasing and equally spaced.
If the last term alone lies in the latest box, each last degree is at
most $2n_{j-1}$, too small. If the last two lie there, the first degree
in each coordinate is at least $n_j-2j>n_{j-1}$. Within one box avoidance
was imposed already. This proves \eqref{eq:glueproduct}. Slicing a finite
product ball by its last coordinate gives $M_{b+1}(n)\le M_b(n)$:
any forbidden triple in a slice lifts by adjoining ratio $1$. Hence
$\Delta(R_b)$ is nonincreasing and has a limit.

The tuples with every component degree in $\{n-1,n\}$, $n\ge1$,
occupy exactly $\prod_i(1-q_i^{-2})$ of $B_n(R_b)$. A surviving nonunit
ratio would have at least one component of positive degree. In that
coordinate its first and last degrees would differ by at least two.
This proves the lower product bound by \eqref{eq:glueproduct}.
Since $0<q_i^{-2}\le1/4$, the infinite product is positive exactly when
the sum of its deficits converges. Together with the cube bound this
completes the dichotomy.

\begin{corollary}[The regime of densities close to one]\label{cor:nearone}
For any sequence of finite products (with no nesting assumption), set
$S=\sum_i q_i^{-2}$ for each product. Along such a sequence,
$\Delta\to0$ exactly when $S\to\infty$, and $\Delta\to1$ exactly when
$S\to0$. In the latter regime the first-order estimate is uniform:
\[
 1-\Delta=S+O(S^{3/2}).
\]
\end{corollary}
\begin{proof}
The upper estimate from word cubes tends to zero when $S\to\infty$
by the variance argument above. The product lower bound is at least
$\exp(-4S/3)$, since $-\log(1-u)\le4u/3$ for $0\le u\le1/4$.
For the other endpoint, slicing any nonempty-dimensional ternary cube
into parallel lines shows $a_k\le2/3$ for $k\ge1$. With
$p_i=3/(q_i^2+q_i+1)$ this gives
\[
 \frac{1-\prod_i(1-p_i)}{3}\le 1-\Delta\le S.
\]
As $\sum_i p_i\ge12S/7$, the left side stays positive when $S$ stays
bounded below. As $S\to0$ each $q_i^{-1}\le\sqrt S$. Thus
$\sum_i p_i=3S+O(S^{3/2})$.
The first two inclusion--exclusion bounds give
$1-\prod_i(1-p_i)=\sum_i p_i+O((\sum_i p_i)^2)$. This proves the
uniform expansion.
\end{proof}

For a finite illustration, the bounds for two copies of $\mathbb F_2$ are
$9/16\le\Delta\le38/49$; for
$\mathbb F_2\times\mathbb F_3\times\mathbb F_5$ they are
$16/25\le\Delta\le322/403$. These use the values
$a_1=a_2=a_3=2/3$ from
\href{https://arxiv.org/abs/1002.0374}{D.~H.~J. Polymath's finite cube calculation}.
The program \texttt{check\_products.py} verifies the chain counts by direct
enumeration at small cutoffs, checks all $37$ lines in a three-coordinate
cell, and computes the displayed bounds in rational arithmetic.

\subsection{Using the repository's Behrend lemma for a lower bound}
\label{sec:behrend-products}
The source lemma
\href{https://github.com/lsernau/grh-lean/blob/72fead102f663d5c7a407964a13dc0cf9d977207/grh/src/New/Reduction/Data/Existence.lean\#L77-L113}{\texttt{exists\_three\_term\_free}}
works with points on a sphere in a digit box. It supplies $S\subset[0,o)$
of size at least $H^d/(dH^2)$ for $(LH)^d\le o$, avoiding nontrivial
weighted three-term congruences with positive weights of sum at most $L$.
In its application to infinite places, multiplicative relations impose
additive constraints on exponents. Here the transfer uses ordinary
midpoints ($L=2$) and projects each polynomial tuple to a single integer.

In particular there is an absolute $C_0$ such that for each $m\ge2$ some
$S\subset\{0,\ldots,m-1\}$ contains no nontrivial three-term arithmetic
progression and has size at least $m\exp(-C_0\sqrt{\log(2+m)})$.
For large $m$ take $d=\lfloor\sqrt{\log_2m}\rfloor$,
$H=\lfloor m^{1/d}/2\rfloor$, and $o=m$. Then $(2H)^d\le m$,
$H^d\ge m/4^d$, and $H^2\le m^{2/d}$. Substitution gives the bound;
a larger absolute $C_0$ covers the finitely many smaller values.

For a random tuple at cutoff $n$ with nonzero components consider its
total degree deficit $D_n=\sum_{i=1}^b(n-\deg f_i)$. For fixed $b$ its
limiting law is the sum $D$ of independent geometric variables with
\[
 \Pr(D_i=t)=(1-q_i^{-1})q_i^{-t}\quad(t\ge0),\qquad
 V=\mathrm{Var}(D)=\sum_{i=1}^b\frac{q_i}{(q_i-1)^2}\le2b .
\]
Indeed the stated geometric probability for a given $t\le n$ is the
exact unconditioned probability of degree $n-t$; zero components have
vanishing probability. An integer interval $I$ of length (cardinality)
$m\le4\sqrt V+3$, enlarged to at least two points if necessary, captures
at least $3/4$ of $D$ by Chebyshev's inequality. As $h$ ranges over
$I-\{0,\ldots,m-1\}$, an interval of $2m-1$ integers, each $t\in I$
belongs to exactly $|S|$ sets $h+S$. Consequently for some $h$,
\[
 \Pr(D\in h+S)\ge\frac{3|S|}{4(2m-1)}.
\]
At cutoff $n$ take the nonzero-component tuples with $D_n\in h+S$.
Along a candidate progression the total degree rises by $\sum_i\deg r_i>0$,
so its three deficits would form a forbidden arithmetic progression.
This admissible finite construction has limiting proportion at least the
last displayed bound. Applying \eqref{eq:glueproduct} and $V\le2b$ proves
the stated lower estimate, after enlarging the absolute constant. In fact
the proof gives $\exp(-C\sqrt{\log(2+V)})$. This application uses the
finite Behrend construction inside the Lean proof and the density
Hales--Jewett theorem from outside it; neither estimate requires RH.

\paragraph{Comparison with a local algebra.}
For even $q$ let $A_b=\mathbb F_q\oplus V$ with $\dim_{\mathbb F_q}V=b-1$
and $V^2=0$. The earlier
\href{https://github.com/VladimirReshetnikov/ProveIt/blob/5ed40f562710fcc2770bbb463dba528a6727f7ae/docs/incoming/grh_method_transfers_research.zip}{ProveIt report},
in its section on finite local rings, proves
$1-Q^{-2}\le\Delta(A_b)\le1-(Q^2+Q+1)^{-1}$ for $Q=q^b$.
One can see the limit $1$ already from the earlier report's
squarefree-gcd construction: in this algebra $(a+v)^2=a^2$, so a
nonunit square producing a nonzero endpoint forces every scalar component
of that endpoint to have a common squared prime divisor. By excluding
all such endpoints one retains natural density $1-q/Q^2$.
For the same $q,b$, the reduced algebra $\mathbb F_q^b$ has the same
cardinality $Q$, yet its maximal upper density tends to zero. Its
coordinate idempotents permit the independent ratios that realize the
ternary word cube. Question~6.5 includes both examples; its general
finite-ring form and the exact finite-factor extrema remain open here.

\section{Further research questions}
\begin{enumerate}[leftmargin=1.5em]
\item Is there an effective algorithm, for all rank-two integral bouquet
voltages, that computes both the eventual polynomial and its least starting
level? The examples suggest reducing the symbol in characteristic $p$,
then handling finitely many exceptional directions over a small cyclotomic
field. An algorithm needs a termination proof and an effective bound on
the precision in those directions.
\item For $p>2m$ in dimension two, Theorem~\ref{thm:moments} reduces the
extremal problem to the orders of binary moment forms. Determine the largest
$A$ when the quadratic form vanishes identically. In higher dimension,
can one express the distribution of $A$ for random step lists through
point counts of the successive moment varieties?
\item For fixed $q$, determine the rate of decay of
$\Delta(\mathbb F_q^b)$. The density Hales--Jewett reduction gives
vanishing and the Behrend construction gives a subexponential-in-$\log b$
lower bound. Can additional polynomial ratios beyond those with entries
$1,x$ substantially sharpen the upper bound?
\item Classify growing products of finite local rings when their radicals
are nonzero. The reduced-ring proof uses coordinate degrees in integral
domains for the lower bound. Nilpotent leading coefficients can change
degree under multiplication, as the square-zero comparison illustrates.
\end{enumerate}

\clearpage
\section*{Verification and reproducibility}
\addcontentsline{toc}{section}{Verification and reproducibility}
The files in \texttt{verification/} use the Python standard library. The
primary graph check constructs each character's Laurent polynomial with
integer coefficients, reduces it modulo the appropriate cyclotomic
polynomial, and changes to the basis $1,(Z-1),\ldots,(Z-1)^{e-1}$. It
does not assume the case distinctions in the proofs. On small graphs
a second check constructs the reduced Laplacian directly and evaluates
its determinant by fraction-free elimination. This last check also tests
the parallel-edge convention at the bottom of the dyadic towers.
\begin{center}\small
\begin{tabular}{lp{.61\linewidth}}
\toprule
Program & Checks performed\\
\midrule
\texttt{check\_graphs.py} & All published DuBose--Vallières numerical levels, cofactors on graphs with at most 90 vertices in these cases, and six examples of the moment formula.\\
\texttt{check\_kataoka.py} & All nine Table 7 rows at $n=1,\ldots,9$; direct cofactors at levels 1--3.\\
\texttt{check\_dv3.py}, \texttt{check\_odd\_rows.py} & Higher-conductor tests through $n=12$ for the unequal-slope dyadic example and Kataoka's indices 5 and 7. The latter also checks the exceptional root-of-unity expansions.\\
\texttt{check\_products.py} & Counts disjoint polynomial chains, tests all 37 lines of a ternary cube of dimension three, and computes finite-factor density bounds exactly.\\
\bottomrule
\end{tabular}
\end{center}
The general statements about all levels follow from the proofs in the
text. The scripts provide both finite base cases and independent tests.
\texttt{source-audit.txt} gives pinned Lean paths, declaration lines, and
file hashes. \texttt{incoming-inventory.txt} records the 41 comparison
archives at the ProveIt revision inspected. The separate local-field note
in this package treats valuation losses over $\mathbb Q_3(\sqrt[3]{3})$.
\begin{thebibliography}{99}\small
\bibitem{Alpoge} L. Alpöge, \emph{Remarks on the Riemann hypothesis},
2026. \url{https://alpo.ge/remark.pdf}.
\bibitem{LeanRepo} L. Sernau and collaborators, \emph{GRH for Dirichlet
$L$-functions, in Lean 4}, 2026. Source, reading guide and verification
notes. \href{https://github.com/lsernau/grh-lean/tree/72fead102f663d5c7a407964a13dc0cf9d977207/grh}{Pinned readable edition}.
\bibitem{DV} S. DuBose and D. Vallières, \emph{On $\mathbb Z_\ell^d$-towers
of graphs}, Algebraic Combinatorics \textbf{6} (2023), 1331--1346.
Theorem 6.2 and Section 7. \url{https://doi.org/10.5802/alco.304}.
\bibitem{Kat} T. Kataoka, \emph{An Iwasawa-type asymptotic formula for
multiple $\mathbb Z_p$-coverings of graphs}, arXiv:2606.03579v1 (2026).
Section 6.3, Tables 6--7. \url{https://arxiv.org/abs/2606.03579}.
\bibitem{DLRV} C. Dion, A. Lei, A. Ray and D. Vallières, \emph{On the
distribution of Iwasawa invariants associated to multigraphs}, Nagoya
Math. J. \textbf{253} (2024), 48--90. \url{https://doi.org/10.1017/nmj.2023.18}.
\bibitem{KM} S. Kleine and K. Müller, \emph{On the growth of the Jacobians
in $\mathbb Z_p^l$-voltage covers of graphs}, Algebraic Combinatorics
\textbf{7} (2024), 1011--1038. \url{https://doi.org/10.5802/alco.366}.
\bibitem{Asada} M. Asada, E. Fourakis, S. Manski, N. McNew, S. J. Miller
and G. Moreland, \emph{Subsets of $\mathbb F_q[x]$ free of 3-term
geometric progressions}, Finite Fields Appl. \textbf{44} (2017), 135--147.
Question 6.5. \url{https://doi.org/10.1016/j.ffa.2016.10.002}.
\bibitem{Behrend} F. A. Behrend, \emph{On sets of integers which contain
no three terms in arithmetical progression}, Proc. Natl. Acad. Sci. USA
\textbf{32} (1946), 331--332. \url{https://doi.org/10.1073/pnas.32.12.331}.
\bibitem{DHJ} D. H. J. Polymath, \emph{A new proof of the density
Hales--Jewett theorem}, Ann. of Math. \textbf{175} (2012), 1283--1327.
\url{https://doi.org/10.4007/annals.2012.175.3.6}.
\bibitem{Moser} D. H. J. Polymath, \emph{Density Hales--Jewett and
Moser numbers}, arXiv:1002.0374. \url{https://arxiv.org/abs/1002.0374}.
\end{thebibliography}

\end{document}
